\begin{abstract}
Dynamic Thermal Management (DTM) is essential for maximizing the performance of modern heterogeneous multicores, such as ARM big.LITTLE architectures, within strict thermal design power (TDP) constraints. However, conventional thermal models used in OS-level scheduling predominantly rely on oversimplified approximations that fail to capture the thermal inertia of physical heatsinks, leading to premature Dynamic Voltage and Frequency Scaling (DVFS) throttling. In this paper, we extend the gem5 full-system simulator with an instrumentation patch that directly exports three-node Cauer RC temperatures, enabling observation of die, package, and heatsink temperatures from simulation statistics. Building on gem5-extracted power and IPC baselines, we design a high-fidelity 3-Node Cauer thermal model (Die--Package--Heatsink) with distinct time constants ($\tau_{die}$\,=\,5\,s, $\tau_{pkg}$\,=\,10\,s, $\tau_{hs}$\,=\,120\,s), and evaluate a package-aware Phase~6 scheduling policy that prioritizes task migration over DVFS throttling. Co-simulation results show that the 3-node model defers the first thermal intervention by 85 simulated seconds compared to a 2-node baseline, and the migration policy reduces DVFS oscillation events from 13 to 1 while lowering modeled package temperature by 4.5\,°C.
\end{abstract}
